%
% 6.006 problem set 6 solutions template
%
\documentclass[12pt,twoside]{article}

\input{macros-sp20}
\newcommand{\theproblemsetnum}{6}

\title{6.006 Problem Set \theproblemsetnum}

\begin{document}

\handout{Problem Set \theproblemsetnum}

\setlength{\parindent}{0pt}
\medskip\hrulefill\medskip

{\bf Name:} Nicolas Iskandar

\medskip

{\bf Collaborators:} None

\medskip\hrulefill

\begin{problems}

\problem  % Problem 1: Dijkstra Practice

\begin{problemparts}

\problempart % Problem 1a
Relaxation order is unique for both algorithms.

\medskip

\textbf{DAG Relaxation:} $(a,b),\ (a,d),\ (d,b),\ (d,e),\ (b,c),\ (b,e),\ (b,f),\ (e,f),\ (f,c)$

\medskip
\textbf{Dijkstra:} $(a,b),\ (a,d),\ (b,c),\ (b,e),\ (b,f),\ (f,c),\ (e,f),\ (d,b),\ (d,e)$

\problempart % Problem 1b
\begin{center}
\begin{tabular}{c|cccccc}
$v$ & $a$ & $b$ & $c$ & $d$ & $e$ & $f$ \\
\hline
$\delta(a,v)$ & 0 & 2 & 3 & 6 & 5 & 4
\end{tabular}
\end{center}

\end{problemparts}

\problem  % Problem 2: Short Circuits
\textbf{Idea.} Standard Bellman--Ford needs $|V|-1$ rounds of relaxing every
edge, because a shortest simple path can use up to $|V|-1$ edges. Here we are
guaranteed every vertex has some shortest path using at most $k$ edges, so we
only need $k$ rounds instead of $|V|-1$.

\textbf{Key fact.} After round $i$ of Bellman--Ford (relaxing every edge $i$
times), the distance array holds the correct shortest-path distance to each
vertex \emph{restricted to paths of at most $i$ edges}. Round $1$ gives the
best 1-edge paths, round $2$ the best 2-edge paths, etc. (Formally this
corresponds to layer $i$ of the duplicated DAG used to prove Bellman--Ford's
correctness; a shared distance array updated in place implements this without
building the layers explicitly.)

\textbf{Algorithm.}
\begin{enumerate}
  \item Initialize $\delta(s)=0$, $\delta(v)=\infty$ for all other $v$. \quad $O(V)$
  \item Repeat $k$ times: relax every edge once (one pass over $E$). \quad $O(kE)$
  \item Return the distance array.
\end{enumerate}

\textbf{Correctness.} Each round extends "best path using $\le i$ edges" to
"best path using $\le i+1$ edges." Since every vertex has an optimal path with
$\le k$ edges by hypothesis, round $k$ is guaranteed to have found it; further
rounds cannot improve anything.

\textbf{Running time.} $O(V)$ for initialization plus $k$ rounds each costing
$O(E)$ gives $O(V+kE)$ total.

\newpage

\problem  % Problem 3: Dynamite Detonation
\textbf{Why elevation changes break a single DAG.} Detonation lowers the
elevation of every clearing in $C_D$, which can \emph{flip} which direction is downhill on any slope touching one of those clearings. A single graph cannot represent both directions of the same edge simultaneously, so a plain shortest-path computation over one graph is
insufficient.

\textbf{Construction.} Build two copies of the clearing graph:
\begin{itemize}
  \item $G_1$: one node per clearing using original elevations $e_i$; a
  directed edge for each slope pointing toward the lower clearing, weighted by
  slope length (drop non-monotonic/flat slopes).
  \item $G_2$: identical construction, but using post-detonation elevations
  $e_i'$ for $c_i \in C_D$.
  \item One zero-weight edge $D_1 \to D_2$ (detonating: instantaneous, costs
  no distance). No edge goes from $G_2$ back to $G_1$, detonation is
  permanent, so this keeps the combined graph acyclic.
  \item A supernode $T$ with zero-weight edges in from $S_1$ and $S_2$,
  merging "reach $S$ without detonating" and "reach $S$ after detonating"
  into a single target.
\end{itemize}

Run DAG Relaxation from $L_1$ to $T$ on the combined graph $G$; the shortest
path automatically chooses whichever strategy (detonate or not) is cheaper.

\textbf{Running time.} $G_1,G_2$ each have $O(n)$ vertices/edges, so $G$ has
$O(n)$ vertices and edges total. DAG Relaxation runs in $O(V+E)=O(n)$.

\problem  % Problem 4: Conservative Cycles
\textbf{Reframing.} Build a directed weighted graph $G$: one vertex per
location, one edge per transition weighted by its work $w_{ij}$. The field is
conservative iff every cycle in $G$ has total weight exactly $0$.

\textbf{Key idea.} Bellman--Ford detects negative-weight cycles as a
byproduct of computing shortest paths. Repurpose this as the actual
detection tool:
\begin{itemize}
  \item Run Bellman--Ford once on the original weights, catches any
  \emph{negative}-weight cycle.
  \item Negate every edge weight and run Bellman--Ford again, a
  positive-weight cycle in the original graph becomes negative-weight in the
  negated graph (weights, and hence cycle totals, flip sign under negation;
  zero-weight cycles stay zero), so this run catches any \emph{positive}-weight
  cycle.
\end{itemize}
If either run reports a negative cycle, the field is not conservative;
otherwise every cycle has weight $0$ and the field is conservative.

\textbf{Supernode.} Add vertex $s$ with a zero-weight edge to every other
vertex (one-directional, out of $s$ only, adds no new cycles). This
guarantees every vertex, and hence every cycle, is reachable from $s$, so a
single Bellman--Ford run from $s$ suffices to catch a negative cycle anywhere
in the graph.

\textbf{Running time.} $V=O(n)$, $E=O(n)$ after adding $s$ and its edges.
Each Bellman--Ford run costs $O(VE)=O(n^2)$; two runs are still $O(n^2)$.

\newpage

\problem  % Problem 5: SparkPlug Derby
\textbf{Why plain Dijkstra fails.} Whether a road is usable depends on
current gas level (uphill travel costs gas and the tank must stay strictly
positive), so a single vertex per intersection cannot capture enough state:
"at $A$ with 5 gas" and "at $A$ with 1 gas" allow different future moves.

\textbf{State-space construction.} Make one vertex $v_i$ per (intersection
$v$, gas level $i$), for $i=1,\dots,g$. For each road $\{u,v\}$ with driving
time $t(u,v)$ and $e(u)\le e(v)$:
\begin{itemize}
  \item \emph{Downhill/flat} $(v_i \to u_i)$ for every $i$, weight $t(u,v)$
  (free in gas).
  \item \emph{Uphill}, costing $\ell=e(v)-e(u)$ gas: edge $(u_i \to v_{i-\ell})$
  only for $i>\ell$, weight $t(u,v)$, this is exactly how "tank stays
  strictly positive" is enforced, by omitting edges that would violate it.
  \item \emph{Refuel} at a station $v$: edge $(v_i \to v_{i+1})$, weight
  $t_G$, for $i<g$. Refueling by $k$ units in one action costs $k\cdot t_G$,
  identical to taking $k$ such edges in a row, so no information is lost by
  only modeling unit refills.
\end{itemize}
All weights are non-negative, and the graph may contain cycles, so run
Dijkstra (not DAG Relaxation) once from $s_g$ (start full). The answer is
$\min_{i} \delta(s_g, t_i)$.

\textbf{Running time.} Average degree $<5 \Rightarrow r=O(n)$; since $g<n$,
the blown-up graph has $O(n^2)$ vertices and edges, built in $O(n^2)$ time.
Dijkstra with a binary heap costs $O(E\log V)=O(n^2\log n)$, dominating the
final $O(n)$-time scan over the $t_i$. Total: $O(n^2\log n)$.

\end{problems}

\end{document}
